We compared the theorem-independent bitset solver (Solver B) with the structure-aware solver (Solver C) on the complete census of 4,681 non-isomorphic connected simple cubic graphs of orders $4,6,\ldots,16$.  The census was generated with nauty 2.9.3 and its layer sizes were independently checked as $1,2,5,19,85,509,$ and $4060$.  Both solvers returned the same optimum on every graph, and every reported witness passed a definition-level maximal acyclic matching checker.  Whole-layer timings used five repetitions through order 12 and three repetitions at orders 14 and 16, with alternating solver order and medians reported.  At order 16, Solver B required 94.18 seconds and Solver C 32.46 seconds, an aggregate speedup of $2.90\times$; the median and geometric-mean per-graph speedups were $4.58\times$ and $4.46\times$, respectively.  Solver C was faster on every order-16 graph and reduced the number of complete candidates from 4,832,310 to 40,158 (99.17\%).  Across all census layers, aggregate time decreased from 100.70 to 37.83 seconds, with no timeouts or discrepancies.  On a separate equality family, Solver B approached the 300-second limit at order 36 and timed out at orders 46 and 56, whereas Solver C solved those two instances in 0.021 and 0.030 seconds.  Since Solver C uses the proved cubic bound, these experiments assess computational performance and are not independent evidence for the theorem.
